%! Orthogonality of the Four Subspaces — Unit 4, Lesson 4.1 — Student Packet Cover Sheet
\documentclass[10pt]{article}
\usepackage{linalg-article}
\usepackage{linalg-boxes}
\usepackage{ltablex}
\keepXColumns

\begin{document}

% Full-bleed burgundy banner
\begin{tikzpicture}[remember picture, overlay]
  \fill[burgundy] (current page.north west)
    rectangle ([yshift=-0.9in]current page.north east);
\end{tikzpicture}
\vspace{-0.8in}

\begin{center}
  {\color{white}\textbf{\LARGE Linear Algebra}}\\[4pt]
  {\color{blushmid}\large\textbf{Unit 4 \quad Orthogonality}}\\[6pt]
  {\color{white}\textbf{\large Lesson 4.1 \quad Orthogonality of the Four Subspaces}}
\end{center}

\vspace{0.22in}
\namedateperiod
\vspace{0.18in}

\begin{learningtargetbox}
\small
\begin{enumerate}[leftmargin=1.5em, itemsep=5pt]
  \item \ldots{} say what it means for two \textbf{subspaces} to be \textbf{orthogonal}, and test it with dot products.
  \item \ldots{} explain \emph{why} the \textbf{row space} $\perp$ the \textbf{nullspace} and the \textbf{column space} $\perp$ the \textbf{left nullspace}.
  \item \ldots{} recognize each pair as an \textbf{orthogonal complement} --- perpendicular \emph{and} filling the whole space ($r+(n-r)=n$).
\end{enumerate}
\end{learningtargetbox}

\vspace{0.14in}

\begin{tocbox}
\small
\renewcommand{\arraystretch}{1.5}
\begin{tabularx}{\linewidth}{@{} c l X r @{}}
\rowcolor{blush}
\textbf{\#} & \textbf{Component} & \textbf{Description} & \textbf{Score} \\
\midrule
1 & Warm-Up        & Spiral review: perpendicular test, nullspace direction, counting dimensions & \blank{1.2cm} \\
2 & Guided Notes   & Orthogonal subspaces, the two perpendicular pairs, and orthogonal complements & \blank{1.2cm} \\
3 & Group Activity & Two Rooms, Two Right Angles --- verifying both pairs and building a complement & \blank{1.2cm} \\
4 & Exit Ticket    & Quick check on the pairing, dimension counts, and complements               & \blank{1.2cm} \\
5 & Homework       & Orthogonal subspaces, both rooms of a matrix, and a complement extension     & \blank{1.2cm} \\
\midrule
  & \multicolumn{2}{r}{\textbf{Total}} & \blank{1.2cm} \\
\end{tabularx}
\end{tocbox}

\vspace{0.10in}

\begin{remindbox}
\small One idea drives the lesson: reading $A\mathbf{x}=\mathbf{0}$ \emph{one row at a time} shows every
nullspace vector is perpendicular to every row --- so the \textbf{nullspace} $\perp$ the \textbf{row
space}. Add the dimension count $r+(n-r)=n$ and the two are \textbf{orthogonal complements}: perpendicular
\emph{and} filling the space. The same story, transposed, pairs the \textbf{column space} with the
\textbf{left nullspace}.
\end{remindbox}

\end{document}
